\chapter{Divisibility}
\label{chapter:divisibility}

\section{The Condition of Divisibility}

The simplest Diophantine equation is the equation $ax=b$ in one variable $x$. The condition of its solvability leads to the concept of \emph{divisibility}.

\begin{definition}
  We say that an integer $a$ \define{divides} an integer $b$ if there exists an integer $x$ such that the equation
  \begin{equation*}
    ax=b
  \end{equation*}
  holds. When $a$ divides $b$, we write $a\mid b$ and we say that $a$ is a \define{divisor} of $b$. Sometimes we also refer to $b$ as the \define{dividend}. When $a$ does not divide $b$, we write $a\nmid b$. An integer $a$ is called \define{even} if $2\mid a$ and is called \define{odd} if $2\nmid a$.
\end{definition}

\begin{example}
  The integers $-1$, $0$, and $1$ have special properties with respect to divisibility:
  \begin{enumerate}
  \item The integers $1$ and $-1$ divide any integer $b$. Indeed, the integer $b$ itself is a solution to the equation $1x=b$, and the integer $-b$ is a solution to the equation $-1x=b$.
  \item Every integer $a$ divides $0$. Indeed, $x=0$ is a solution to the equation $ax=0$.
  \item If $0\mid b$ for some integer $b$, then it follows that $b=0$. Indeed, if $0x=b$ has a solution, then $b=0$ because $0x=0$. Furthermore, we observe that \emph{every} integer is a solution to the equation $0x=0$. It follows that $0\nmid 1$, and indeed that $0\nmid b$ for any nonzero integer $b$.
  \item Similarly, if $a\mid 1$ for some integer $a$, then it follows that $a=\pm 1$. Indeed, if $ax=1$ has a solution, then the absolute value $|a|$ of $a$ cannot be greater than $1$, because its nonzero multiples would be of absolute value greater than $1$. Therefore it follows that $|a|=1$, which proves the claim. In the situation where $a\mid 1$, we say that $a$ is \define{invertible}, or that it is a \define{multiplicative unit}, because the solution to the equation $ax=1$ is a multiplicative inverse of $a$.
  \end{enumerate}
\end{example}

%If $a$ is a nonzero integer, then there is always at most one solution to the equation $ax=b$. In other words, if a nonzero integer $a$ divides an integer $b$, then its quotient is uniquely defined. On the other hand, we just saw that every integer is a solution to the equation $0x=0$. The quotient is therefore not uniquely defined in this case. In the following lemma we remove this ambiguity by imposing an extra bound on the solution, that $|x|\leq |b|$. In the case where $a$ is nonzero, then this condition is automatically satisfied, while in the case where $a=0$ it imposes the zero solution.

\begin{lemma}
  Suppose that $a\mid b$, where $a$ is a nonzero integer. Then there is a unique integer $x$ satisfying the equation $ax=b$. We will write $b/a$ for this unique integer, which is called the \define{quotient} of $b$ divided by $a$.
\end{lemma}

\begin{proof}
%  There are two cases to consider: either $a=0$ or $a\neq 0$. In the case where $a=0$, having a solution $ax=b$ implies that $b=0$. In this case, any $x$ is a solution to the equation $ax=b$. However, only one solution $x$ satisfies $|x|\leq b$, namely $x=0$. Thus, the quotient $0/0$ is set to be $0$.
%
  By the assumption that $a\neq 0$, it follows that the function $x\mapsto ax$ is \emph{injective} {\color{Red} [needs reference or proof]}. This means that $ax=ay$ implies $x=y$. Consequently, there can be at most one integer $x$ such that $ax=b$.
\end{proof}

\begin{remark}
  While we have given a definition of divisibility in natural language, the formal definition of divisibility features the \emph{existential quantifier} $\exists$. When we write the logical formula $\exists_{(x\in X)}P(x)$, this means that there exists an element $x\in X$ for which $P(x)$ holds. The definition of divisibility translates to the following logical formula:
  \begin{equation*}
    a\mid b := \exists_{(x\in\mathbb{Z})}ax=b.
  \end{equation*}
  In order to use existential quantification in proofs, it is useful to know how to break them down systematically. There are two cases to consider: (1) How to \emph{prove} an existence claim, and (2) how to \emph{use} an existence claim as an assumption to prove something else.
  \begin{enumerate}
  \item The principal way to prove that there exists an element $x\in X$ such that $P(x)$ holds, is to construct an element $x$ of $X$ and prove that the property $P(x)$ holds.
  \item To use that there exists an $x\in X$ such that $P(x)$ holds in the proof of another property $Q$, you may assume to have an element $x$ and you may assume that the property $P(x)$ holds. However, in this scenario we may not assume anything else about $x$ unless that is otherwise warranted.
  \end{enumerate}
  In the case of divisibility, in order to prove that $a\mid b$ holds we must find a solution to the equation $ax=b$, and in order to use that $a\mid b$ holds as an assumption, we may assume $x$ is an integer for which the equation $ax=b$ holds.
\end{remark}

\begin{proposition}
  If $d$ divides $a$ and $b$, then $d$ divides $ax+by$ for any two integers $x$ and $y$.
\end{proposition}

\begin{proof}
  Suppose that $d$ divides both $a$ and $b$. Let $u$ and $v$ be integers such that $du=a$ and $dv=b$. Then we have
  \begin{equation*}
    ax+by = dux + dvy = d(ux + vy),
  \end{equation*}
  which shows that $ax + by$ is a multiple of $d$.
\end{proof}

Sometimes, the hardest thing is to figure out how to rigorously prove a property that looks completely self-evident. The following proposition, for instance, asserts that in any set of $k$ consecutive integers, exactly one of them is divisible by $k$.

There are, as is commonly the case in mathematics, several ways to approach this problem, but the question is how to generate ideas towards a full proof. One good way of generating ideas is by looking at special cases. In our problem at hand, a special case of interest is the set $\{0,1,\ldots,k-1\}$. There should be exactly one element in the set $\{0,1,\ldots,k-1\}$ divisible by $k$.

Indeed, this is a direct consequence of the Euclidean division theorem: The only element in this set that has remainder $0$ after division by $k$ is the integer $0$. This suggests using the Euclidean division theorem for the general claim as well. Furthermore, we will also use the common proof technique of assuming something \emph{without loss of generality}.

\begin{proposition}
  In any set of $k>0$ consecutive integers
  \begin{equation*}
    \{a,a+1,\ldots,a+k-1\}
  \end{equation*}
  there is exactly one element divisible by $k$.
\end{proposition}

\begin{proof}
  First, consider the function
  \begin{equation*}
    r : \{a,a+1,\ldots,a+k-1\} \to \{0,1,\ldots,k-1\},
  \end{equation*}
  where $r(x)$ is defined to be the remainder of $x$ after division by $k$. We claim that this map is injective, meaning that if $r(x)=r(y)$ for two elements $x$ and $y$, then $x=y$. To see this, consider two elements $x$ and $y$ in the set $\{a,\ldots,a+k-1\}$. Without loss of generality we may assume that $x\leq y$. Indeed, either $x\leq y$ or $y\leq x$ holds, and these situations are completely similar, so nothing is lost if we just assume that $x\leq y$.

  We have the equalities $x=q(x)k+r(x)$ and $y=q(y)k+r(y)$, where $q(x)$ and $q(y)$ are the quotients of $x$ and $y$ after division by $k$, respectively. Since we assumed that $r(x)=r(y)$, it follows that $y-x=(q(y)-q(x))k$. In other words, $y-x$ is divisible by $k$. Furthermore, we have $0\leq y-x <k$ because both $x$ and $y$ are in the set $\{a,\ldots,a+k-1\}$. It follows from Euclidean division theorem that the only integer $0\leq z <k$ divisible by $k$ is $0$, so we find that $y-x=0$. We conclude that $x=y$.

  It now follows that every element $y\in\{0,1,\ldots,k-1\}$ is the value of $r$ of exactly one element $x\in\{a,a+1,\ldots,a+k-1\}$, because both sets have $k$ elements and the map $r$ doesn't take any value more than once. In particular, there is exactly one element $x\in\{a,a+1,\ldots,a+k-1\}$ whose remainder after division by $k$ is $0$. In other words, exactly one element of this set is divisible by $k$.
\end{proof}

\section{The Ordering by Divisibility}
The divisibility relation equips the sets of integers and natural numbers with a useful extra structure, a \emph{partial ordering} of their elements. The set of natural numbers equipped with the divisibility relation is a poset, while the set of integers equipped with the divisibility relation is a preorder, a slightly weaker structure.

\begin{definition}
  A \define{preorder} consists of a set $X$ and an ordering $\leq$ of the elements of $X$, such that
  \begin{enumerate}
  \item The ordering is \define{reflexive}. This means that $x\leq x$ holds for every $x$.
  \item The ordering is \define{transitive}. This means that if $x \leq y$ and $y \leq z$ both hold, then $x\leq z$ also holds.
  \end{enumerate}
  A \define{poset}, or \define{partially ordered set} in full, is a preorder $(X,{\leq})$ satisfying additionally the condition:
  \begin{enumerate}
  \setcounter{enumi}{2}
  \item The ordering is \define{antisymmetric}. This means that if both $x\leq y$ and $y\leq x$ hold, then $x=y$.
  \end{enumerate}
\end{definition}

\begin{theorem}
  The set of integers equipped with the divisibility relation $\mid$ is a preorder, and the set of natural numbers equipped with the divisibility relation $\mid$ is a poset.
\end{theorem}

\begin{proof}
  We first show that the set of integers with divisibility is a preorder. To see that the divisibility relation is reflexive, note that $x=1$ is a solution of the equation
  \begin{equation*}
    ax=a.
  \end{equation*}
  Therefore, it follows that $a\mid a$ for every integer $a$. To see that the divisibility relation is transitive, assume that $a\mid b$ and $b\mid c$. That is, assume that we have integers $x$ and $y$ such that $ax=b$ and $by=c$. Then it follows that
  \begin{equation*}
    a(xy)=(ax)y=by=c.
  \end{equation*}
  In other words, $c$ is divisible by $a$. This completes the proof that $\mathbb{Z}$ equipped with the divisibility relation is a preorder.

  To see that the natural numbers equipped with divisibility forms a poset, it remains to show that the divisibility relation is antisymmetric on the natural numbers. Suppose that $m\mid n$ and $n\mid m$; that is, that there are natural numbers $x$ and $y$ so that $mx=n$ and $ny=m$. By transitivity of divisibility it follows that $m(xy)=m$.
\end{proof}

\begin{figure}[t]
  \centering
  \begin{tikzpicture}[
    y=1.2cm,
    draw=CornflowerBlue,
    thick,
    hasse node/.style={
      draw,
      circle,
      minimum size=7mm,
      inner sep=0pt,
      fill=white
    },
    every node/.style={hasse node},
    overpass/.style={
      preaction={
        draw=white,
        line width=3pt,
        shorten <=2pt,
        shorten >=2pt
      }
    },
    alt={
      The divisors of the number 120 are displayed in six levels.
      The first level contains only the number 1;
      the second level contains the numbers 2, 3, and 5;
      the third level contains the numbers 4, 6, 10, and 15; 
      the fourth level contains the numbers 8, 12, 20, and 30;
      the fifth level contains the numbers 24, 40, and 60;
      and the top level contains only the number 120.
      Two numbers in successive levels are connected by an edge if the lower number divides the higher number.
      }
    ]
    \node (1)  at (0,0) {$1$};
    
    \node (2)  at (-2,1) {$2$};
    \node[draw=CornflowerBlue!50!Gray,draw opacity=.7] (3)  at (0,1) {$3$};
    \node (5)  at (2,1) {$5$};
  
    \node (4)  at (-3,2) {$4$};
    \node[draw=CornflowerBlue!50!Gray,draw opacity=.7] (6)  at (-1,2) {$6$};
    \node (10)  at (1,2) {$10$};
    \node (15) at (3,2) {$15$};
  
    \node (8)  at (-3,3) {$8$};
    \node[draw=CornflowerBlue!50!Gray,draw opacity=.7] (12) at (-1,3) {$12$};
    \node (20) at (1,3) {$20$};
    \node (30) at (3,3) {$30$};

    \node (24) at (-2,4) {$24$};
    \node (40) at (0,4) {$40$};
    \node (60) at (2,4) {$60$};

    \node (120) at (0,5) {$120$};

    \draw (1) -- (2);
    \draw (1) -- (5);

    \draw (5) -- (15);
    \draw[overpass] (5) -- (10);
    \draw (2) -- (4);
    \draw[overpass] (2) -- (10);
      
    \draw (15) -- (30);
    \draw (10) -- (30);
    \draw[overpass] (10) -- (20);
    \draw (4) -- (8);
    \draw[overpass] (4) -- (20);
    
    \draw (30) -- (60);
    \draw (20) -- (60);
    \draw[overpass] (20) -- (40);
    \draw (8) -- (24);
    \draw[overpass] (8) -- (40);
    
    \draw (24) -- (120);
    \draw (40) -- (120);
    \draw (60) -- (120);

    \begin{scope}[
      on background layer={
        draw=CornflowerBlue!50!Gray,
        draw opacity=.7,
        thick}
      ]

      \draw (1) -- (3);
      
      \draw (3) -- (6);
      \draw (3) -- (15);
      \draw (2) -- (6);
      
      \draw (6) -- (12);
      \draw (6) -- (30);
      \draw (4) -- (12);
      
      \draw (12) -- (24);
      \draw (12) -- (60);
    \end{scope}
  \end{tikzpicture}
  \caption{The Hasse diagram of the divisors of $120$.}
  \label{figure:hasse-diagram-120}
\end{figure}

The diagram displayed in \zcref{figure:hasse-diagram-120} is the \emph{Hasse diagram} of the positive divisors of the number $120$, a fragment of the poset of the natural numbers ordered by divisibility.

Despite lacking antisymmetry, the preorder of the integers with divisibility is essentially equivalent to the poset of natural numbers with divisibility. For every integer $a$ there is a unique natural number $n$ satisfying $n\mid a$ and $a\mid n$. This unique number is of course the absolute value of $a$. The absolute value function $x\mapsto |x|$ preserves divisibility, which means that $a\mid b$ implies $|a| \mid |b|$, and it satisfies the logical equivalence
\begin{equation*}
  (a\mid b)\land (b\mid a) \Leftrightarrow |a|=|b|
\end{equation*}
for any integer $a$ and any natural number $n$. Order preserving maps $f : (P,\leq_P) \to (Q,\leq_Q)$ from a preorder $(P,\leq_P)$ to a poset $(Q,\leq_Q)$ that satisfy the logical equivalence $(x\leq_P y)\land (y\leq_P x) \Leftrightarrow f(x)=f(y)$ for every $x,y\in P$ are also called \emph{poset reflections}. Every preorder has a poset refection, but we will not need to go into that aspect of order theory in this course.

\section{Ideals of Integers}

A fruitful way of studying the divisibility relation is by studying the set of \emph{all} integers divisible by a given integer. That is, instead of focusing on one instance of divisibility, we turn our attention to the totality of all integers divisible by a given integer, and we ask questions such as: What properties does the set of all integers divisible by an integer $a$ have? This abstraction leads us to the concept of \emph{ideal}, which was first defined and studied by Richard Dedekind \cite{dedekind1871,kleiner1998}.

\begin{definition}
  An \define{ideal of integers} is a subset $I\subseteq \mathbb{Z}$ that contains $0$, and contains the linear combination $kx+ly$ for any $x,y\in I$, and any $k,l\in\mathbb{Z}$.  
\end{definition}

\begin{example}
  The subset $\{0\}\subseteq{\mathbb{Z}}$ satisfies the conditions of being an ideal in a trivial way; this ideal is called the \define{zero ideal}.
\end{example}

\begin{example}
  The subset $(a):=\{ka \mid k\in\mathbb{Z}\}$ also satisfies the conditions of being an ideal; ideals of this form are called \define{principal ideals}, since they are generated by a single integer.  
\end{example}

In the following theorem we show that every ideal of integers is a principal ideal. By virtue of this properties, we also call $\mathbb{Z}$ a \define{principal ideal domain}.

\begin{theorem}
  \label{theorem:integers-principal-ideal-domain}
  For every ideal $I\subseteq\mathbb{Z}$ there is a unique natural number $n$ such that $I=(n)$.
\end{theorem}

\begin{proof}
  There are two cases to consider: Either $I$ is the zero ideal, or $I$ contains a nonzero integer. In the first case it is clear that $I=(0)$, and that $0$ is the unique natural number $n$ such that $I=(n)$. Thus we focus on the nonzero case.

  Suppose that $I$ contains a nonzero integer $x$. Then $I$ contains the integer $|x|$, because $|x|=\pm x$. In particular, $I$ contains a positive integer. Therefore it follows that $I$ contains a least positive integer $n$. To see that $I=(n)$, consider an element $x\in I$. By the Euclidean division theorem, it follows that $x=qn+r$, where $0\leq r<n$. Notice that $r\in I$, since it is the difference $x-qn$ of two elements in $I$. Since $r$ is nonnegative and $n$ is the least positive integer in $I$, it follows that $r=0$, which gives us that $x=qn$. Thus, every element in $I$ is a multiple of $n$.

  Now suppose that $I=(n)$ and $I=(m)$ for two nonzero natural numbers $n$ and $m$. Then it follows that $m=qn$ and $n=pm$. Thus, we see that $m=(pq)m$, which implies that $pq=1$. Furthermore, both $p$ and $q$ are positive integers, so we have $p=1$ and $q=1$. This shows that $m=1n=n$, establishing uniqueness.
\end{proof}

The set of ideals of $\mathbb{Z}$ is ordered by inclusion: We write $I\subseteq J$ for two ideals $I$ and $J$, if every element of $I$ is an element of $J$. Note that $(b)\subseteq (a)$ holds if and only if $a\mid b$. The theory of ideals of integers is therefore closely related to the theory of divisibility of integers.

\exercises

\begin{starter}
  \exercise Draw the Hasse diagrams for the posets of positive divisors of the following numbers: $360$, $900$, $1024$, $1296$, and $2310$.
  \exercise Show that $a^n-1$ is divisible by $a-1$ for every integer $a$.
  \exercise Show that $a^n+b^n$ is divisible by $a+b$ for any odd natural number $n$.
  \exercise Suppose that $a\mid b$, where $a$ is assumed to be nonzero.
  \begin{subexenum}
  \item Show that $d\mid b/a$ if and only if $ad\mid b$.
  \item Show that for any integer ideal $I$, $b/a\in I$ if and only if $b\in aI$, where $aI=\{ax \mid x\in I\}$.
  \end{subexenum}
  \exercise Suppose that $a\mid b$ and $b\mid c$, where $a\neq 0$.
  \begin{subexenum}
  \item Show that if $d\mid b/a$ then $d\mid c/a$.
  \item Show that for any integer ideal $I$, if $b/a\in I$ then $c/a\in I$.
  \end{subexenum}
  \exercise Show that if $a$ and $ab$ are $n$\supth{} powers, with $a\neq 0$, then so is $b$.
  \exercise Suppose that $a\neq 0$. Show that if $ab$ and $ac$ are squares, then so is $bc$.  
\end{starter}

\begin{routine}
  \exercise 
  \begin{subexenum}
  \item Show that $a^2-b^2$ is divisible by $8$ for any two odd integers $a$ and $b$.
  \item Show that $(a^2-1)/8$ is a triangular number for any odd integer $a$.
  \end{subexenum}
  \exercise
  \label{exercise:div-third-power-minus-first-power}
  \begin{subexenum}
  \item Show that $a^3-a$ is divisible by $6$ for every integer $a$.
  \item Show that
    \begin{equation*}
      \sum_{k=1}^{a-1}\binom{k+1}{2}=\frac{a^3-a}{6}
    \end{equation*}
    for every $a\geq 0$.
  \item Use the previous identity to give a geometric interpretation of the numbers $(a^3-a)/6$.
  \end{subexenum}
  \exercise
  \label{exercise:div-fifth-power-minus-first-power}
  Show that $a^5-a$ is divisible by $30$ for every integer $a$.  
  \exercise Show that for any two positive integers $a \leq b$ there is a unique integer $a\leq x \leq b$ satisfying the logical equivalence
  \begin{equation*}
    2^k\mid x\quad\Leftrightarrow\quad 2^k\mid \lcm(a,\ldots,b)
  \end{equation*}
  for every $k\geq 0$.
\end{routine}

\begin{challenge}
  \exercise
  \label{exercise:addition-formula-fibonacci}
  \begin{subexenum}
  \item Prove the \define{Fibonacci addition formula} $F_{m+n+1} = F_{m+1}F_{n+1} + F_{m}F_{n}$.
  \item For any integer $d$ and any two natural numbers $m$ and $n$, show that if two of the three following conditions hold, then so does the third:
    \begin{enumerate}
    \item $d\mid F_m$.
    \item $d\mid F_n$.
    \item $d\mid F_{n+m}$.
    \end{enumerate}
%  \item Show that the Fibonacci sequence preserves and reflects divisibility: We have $m\mid n$ if and only if $F_m\mid F_n$.
  \end{subexenum}
  \exercise
  \label{exercise:div-product-sum-of-squares}
  Show that if $x$ and $y$ are positive integers such that $xy$ divides $x^2+y^2$, then $x=y$.
  \exercise Show that in any subset of $n+1$ integers from the set $\{1,\ldots,2n\}$, there is a pair of distinct integers such that one divides the other. Does the claim still hold for every subset of $n$ integers from the set $\{1,\ldots,2n\}$?
  \exercise Consider two relatively prime positive integers $m$ and $n$, and define a sequence $(a_k)$ of natural numbers by
  \begin{equation*}
    a_0=m+n \qquad\text{and}\qquad a_{i+1}=a_i^2-na_i+n
  \end{equation*}
  Show that $\gcd(a_i,a_j)=1$ for all $i\neq j$ \cite{mohanty1978}.
  \exercise Show that for $n>1$ there are no solutions in the positive integers to the equation
  \begin{equation*}
    a(a+1)(a+2)=b^n.
  \end{equation*}
  More generally, Paul Erd\H{o}s and John Lewis Selfridge showed that no product of at least two consecutive integers can be an $n$\supth{} power for $n>1$ \cite{erdosselfridge1975}.
  \exercise Find all solutions to the equation
  \begin{equation*}
    3^m-2^n=1.
  \end{equation*}
\end{challenge}

%%% Local Variables:
%%% mode: latex
%%% TeX-master: "number-theory"
%%% TeX-engine: luatex
%%% End:

